% BHU PYQ -- Professional Exam Template
\documentclass[11pt,a4paper]{article}

\usepackage[a4paper,top=2.5cm,bottom=2.5cm,left=2.8cm,right=2.8cm,headheight=14pt]{geometry}
\usepackage{amsmath,amssymb}
\usepackage[T1]{fontenc}
\usepackage[utf8]{inputenc}
\usepackage{lmodern}
\usepackage{microtype}
\usepackage{fancyhdr}
\usepackage{enumitem}
\usepackage{hyperref}
\usepackage{lastpage}
\usepackage{parskip}

\hypersetup{colorlinks=true,urlcolor=black,linkcolor=black,
  pdftitle={B.Sc. (Hons.) Zoology Sem-V ZOB-501 -- BHU PYQ}}

\pagestyle{fancy}
\fancyhf{}
\renewcommand{\headrulewidth}{0.4pt}
\renewcommand{\footrulewidth}{0.4pt}
\fancyhead[L]{\small\textit{Banaras Hindu University}}
\fancyhead[R]{\small\textit{Zoology Paper: ZOB-501}}
\fancyfoot[C]{\small Page \thepage\ of \pageref{LastPage}}

\newlist{parts}{enumerate}{3}
\setlist[parts,1]{label=(\alph*),leftmargin=*,itemsep=4pt,topsep=3pt}
\setlist[parts,2]{label=(\roman*),leftmargin=*,itemsep=2pt,topsep=2pt}
\setlist[parts,3]{label=(\arabic*),leftmargin=*,itemsep=2pt,topsep=2pt}
\newcommand{\pts}[1]{\hfill\textit{[#1]}}

\begin{document}

\begin{center}
  {\large\bfseries BANARAS HINDU UNIVERSITY}\\[2pt]
  {\normalsize B.Sc. (Hons.) Semester V Examination, 2016--17}\\[6pt]
  \rule{0.8\linewidth}{0.5pt}\\[6pt]
  {\large\bfseries Zoology}\\[4pt]
  {\large\bfseries Paper No. ZOB-501 : Functional Anatomy of Non-chordate}\\[4pt]
  {\normalsize Time: 3 hours \hfill Full Marks: 70}
\end{center}

\rule{\linewidth}{0.4pt}

\smallskip
\noindent\textit{Note: Write your Roll No. at the top immediately on the receipt of this question paper. The figures in the right-hand margin indicate marks. Answer six questions including Question 1 which is compulsory.}

\bigskip
\noindent\textbf{Question 1.} Answer the following parts:
\begin{parts}
  \item Define the following terms:\pts{$4 \times 1 = 4$}
  \begin{parts}
    \item Stigma
    \item Scolex
    \item Medusa
    \item Gregarines.
  \end{parts}

  \item State whether the following statements are True or False. Give brief justification to your answer:\pts{$4 \times 2 = 8$}
  \begin{parts}
    \item Gametocyst in \textit{Monocystis} is generated from sporozoite.
    \item Fertilium is produced by male heteronereis.
    \item The correct sequence of larval stages in the development of \textit{Fasciola} is Miracidia, Radiae, Sporocyst, Metacercaria and Cercaria.
    \item Surface tegument of \textit{Fasciola} consists of microtriches for absorption of readymade food in host.
  \end{parts}

  \item Differentiate between the following:\pts{$4 \times 2 = 8$}
  \begin{parts}
    \item Young and Mature trophozoite of \textit{Monocystis}
    \item Bipinnaria and Pluteus
    \item Hydrozoa and Scyphozoa
    \item Aplacophoran and Polyplacophoran foot.
  \end{parts}
\end{parts}

\pagebreak

\medskip\noindent\textbf{Question 2.} Write the name of different types of canal system in sponges with example. Explain the Syconoid type of canal system with a neat diagram.\pts{$2 + 8 = 10$}

\medskip\noindent\textbf{Question 3.} Describe the structure of \textit{Euglena} with a neat diagram. Comment on euglenoid movement.\pts{$7 + 3 = 10$}

\medskip\noindent\textbf{Question 4.} Describe the structure of polyp in \textit{Obelia} with a neat diagram.\pts{10}

\medskip\noindent\textbf{Question 5.} Discuss the structure and significance of Trochophore larva in Annelida.\pts{10}

\medskip\noindent\textbf{Question 6.} Discuss the different types of larva in Echinoderms with labelled sketch.\pts{10}

\medskip\noindent\textbf{Question 7.} Illustrate the life cycle of \textit{Ascaris}. Mention the pathogenesis caused by \textit{Ascaris}.\pts{$8 + 2 = 10$}

\medskip\noindent\textbf{Question 8.} Write short notes on any two of the following in not more than 400 words each:\pts{$2 \times 5 = 10$}
\begin{parts}
  \item Theories of origin of Metazoa
  \item Parasitic adaptation of \textit{Fasciola}
  \item Hydrostatic system of \textit{Asterias}.
\end{parts}

\vfill
\rule{\linewidth}{0.4pt}
\begin{center}\small
\href{https://bhu-exam.vercel.app/}{Click Here}\\[4pt]\href{https://me-aryan.vercel.app/}{Click Here}\\[4pt]\href{https://quashan-me.vercel.app/}{Click Here}
\end{center}

\end{document}
